% Documentary companion volume. Independent source, without automatic compilation.
\documentclass[a4paper,10pt]{article}
\usepackage[top=23mm,bottom=23mm,left=20mm,right=20mm,headheight=15pt]{geometry}
\usepackage{fontspec,amsmath,unicode-math}
\usepackage[english]{babel}
\usepackage{microtype,booktabs,longtable,array,needspace,fancyhdr}
\usepackage[unicode,hidelinks,bookmarksopen=true,bookmarksopenlevel=1,bookmarksnumbered=false]{hyperref}
\setmainfont{STIXTwoText-Regular.otf}[BoldFont=STIXTwoText-Bold.otf,ItalicFont=STIXTwoText-Italic.otf,BoldItalicFont=STIXTwoText-BoldItalic.otf]
\setmathfont{STIXTwoMath-Regular.otf}
\renewcommand{\normalsize}{\fontsize{10.5}{13.4}\selectfont}\normalsize
\setlength{\parindent}{1em}\setlength{\parskip}{3pt}
\setlength{\emergencystretch}{2em}
\clubpenalty=10000\widowpenalty=10000\raggedbottom
\clubpenalties 3 10000 10000 0
\widowpenalties 3 10000 10000 0
\AddToHook{cmd/section/before}{\clearpage}
\AddToHook{cmd/subsection/before}{\Needspace{7\baselineskip}}
\AddToHook{cmd/subsubsection/before}{\Needspace{6\baselineskip}}
\pagestyle{fancy}\fancyhf{}
\fancyhead[L]{Structural and Metrological Catalogue of Particles}
\fancyfoot[C]{\thepage}
\hypersetup{pdftitle={Structural and Metrological Catalogue of Particles},pdfauthor={Oumar Haidara Fall; Rubén Ramos Balsa},pdfsubject={Appendix of States, Routes and Masses}}
\newcommand{\CatalogRouteReference}{Sections~\ref{sec:vi-324-rutas} and~\ref{sec:vi-inventario-metrologico} of this catalogue publish the 324 oriented routes and the 855 metrological records, respectively.}
\begin{document}
\thispagestyle{empty}
\begin{center}
{\fontsize{8}{10}\selectfont Compilation manuscript for academic transfer\par}
\vspace{12mm}
{\fontsize{17}{20.5}\selectfont\bfseries Structural and Metrological Catalogue of Particles\par}
\vspace{6pt}
{\fontsize{11}{13.5}\selectfont Appendix of States, Routes and Masses\par}
\vspace{11pt}
{\fontsize{11.5}{14}\selectfont Oumar Haidara Fall \textperiodcentered\ Rubén Ramos Balsa\par}
\vspace{4pt}
{\fontsize{8}{10}\selectfont Coauthorship without precedence of contribution\par}
\vspace{3pt}
{\fontsize{7.5}{9.5}\selectfont Documentary integration and compilation generated by ChatGPT -- Astra.\par}
\end{center}
\vspace{7pt}
\begin{center}\bfseries Presentation\end{center}
\begingroup
\fontsize{10.5}{12.8}\selectfont
\setlength{\parskip}{2pt plus 1pt minus .5pt}
\noindent
\begin{minipage}{\textwidth}
This catalogue brings together the structural inventory of \(324\)
oriented routes and \(855\) metrological records of masses and widths.
Its organisation makes each route's identity and fields consultable and
the metrological observations locatable, without equating the number of documentary records
with that of individualised predictions. It accompanies \emph{Particles,
Persistence and Derivation of the Mass Spectrum}, which develops the
definitions and proofs constructing the particle, the atlas, and the
mass operator. The complete structural and comparison data domain used
by that development is published here.

The first section distinguishes \(13\) multisections, \(56\) route
families, and \(23\) inventory classes.
Section~\ref{sec:vi-324-rutas} presents the \(324\) routes, preserving
every field of their two concordant records.
Section~\ref{sec:vi-inventario-metrologico} brings together \(471\) masses
and \(384\) widths from the fixed metrological edition, with their
identifiers, statuses, units, and sources. Central values, intervals,
bounds, and widths retain their own meanings. An observation's presence
in the inventory makes it locatable for comparison; its physical
identification requires the reader and chart specified in the article.

Every card preserves its complete record, including provenance fields
and the distinctions between routes, families, and observables. The
linked index, bookmarks by cell and orientation, and repeated identifier
on continuation pages make it possible to traverse the entire set and
return directly to the selected route. Keeping this data appendix
separate maintains legible cards without abbreviating their records.\par
\end{minipage}
\par\endgroup
\clearpage
\tableofcontents
\clearpage
\input{sections/10a_dependencia_activa_l_g.tex}
% Generated by technical/translate_catalogue.py from the frozen catalogue. Do not edit manually.
\clearpage
\section{Catalogues of multisections, route types and sector classes}
\label{sec:vi-catalogos-estructurales}

This appendix organizes three distinct levels of the internal inventory. The
13 multisections group atlas cells; the 56 collections classify internal routes;
the 23 classes document their articulation with physical sectors and
representations. These cardinalities do not count the same object and
cannot be added to obtain a number of particles. The tables accompany the
definitions and proofs of the article: their purpose is to locate each
record, its incidences and its readers, not to replace a proof with a census.

The documentary order of the rows, their identifiers and all repetitions
appearing in the sources are preserved. Literal identifiers allow the tables
to be checked against the JSON files. A dash denotes an empty source cell;
it does not represent zero. The files also preserve every field and digit,
the original statuses and verifiable provenance.
\CatalogRouteReference
% The reference is selected by the article or catalogue entry point.

\subsection{The 13 multisections of the atlas}

For a multisection \(M\), let \(\mathcal C(M)\) be its set of cells and
let \(\gamma_{B_1}(c)\) be the primary route preserved by the atlas in cell
\(c\). In this structural appendix, the fiber is defined by
\[
 B_1(M):=
 \operatorname{span}\bigl\{e_{\gamma_{B_1}(c)}:c\in\mathcal C(M)\bigr\}
 \subseteq \ell^2(\mathcal R_{\mathrm{adm}}).
\]
The primary routes of distinct cells determine distinct basis vectors;
consequently,
\(\operatorname{rank}B_1(M)=\dim B_1(M)=\#\mathcal C(M)\).
This definition fixes the object to which the table's rank refers; the
collections column counts the route configurations present, not physical species.
The integer spectra of the readers \(K_D\) and \(K_\Omega\) are preserved
separately. In the literal notation, \texttt{x} indicates multiplicity.
The 13 rows have the documentary status of an exact diagonal operator on
the fiber \(B_1\). Their representation as a physical scalar requires the
declared transport, except for the rank-one selection provided for in the
record itself. This statement reproduces the recorded scope and does not
turn every spectrum into a unique mass.

\begingroup
\setlength{\tabcolsep}{4pt}
\setlength{\extrarowheight}{4pt}
\renewcommand{\arraystretch}{1.13}
\begin{longtable}{@{}>{\raggedright\arraybackslash}p{\dimexpr 0.19871795\linewidth-2\tabcolsep\relax}>{\raggedright\arraybackslash}p{\dimexpr 0.19871795\linewidth-2\tabcolsep\relax}>{\raggedright\arraybackslash}p{\dimexpr 0.30769231\linewidth-2\tabcolsep\relax}>{\raggedright\arraybackslash}p{\dimexpr 0.29487179\linewidth-2\tabcolsep\relax}@{}}
\caption{The 13 multisections, without row selection.}\label{tab:vi-multisecciones}\\
\toprule
Multisection & Rank and incidence & Classification and depth & Integer readers \\
\midrule
\endfirsthead
\multicolumn{4}{l}{\textit{Continuation}}\\
\toprule
Multisection & Rank and incidence & Classification and depth & Integer readers \\
\midrule
\endhead
\midrule
\multicolumn{4}{r}{Continued on the next page}\\
\endfoot
\bottomrule
\endlastfoot
% VI-CATALOG-ROW multisecciones 1
charged\_\allowbreak{}lepton\_\allowbreak{}G0 & \textit{Rank}: 1\par \textit{Cells (identifiers)}: 41\par \textit{Collections}: 1 & \textit{Collection}: charged\_\allowbreak{}lepton\par \textit{Depth}: G0\_\allowbreak{}center\par \textit{Color}: none & \textit{K D}: (80,\allowbreak{}54,\allowbreak{}6)\allowbreak{}x1\par \textit{K Ω}: (80,\allowbreak{}54,\allowbreak{}6)\allowbreak{}x1 \\
% VI-CATALOG-ROW multisecciones 2
charged\_\allowbreak{}lepton\_\allowbreak{}G2 & \textit{Rank}: 8\par \textit{Cells (identifiers)}: 21,\allowbreak{}23,\allowbreak{}25,\allowbreak{}39,\allowbreak{}43,\allowbreak{}57,\allowbreak{}59,\allowbreak{}61\par \textit{Collections}: 8 & \textit{Collection}: charged\_\allowbreak{}lepton\par \textit{Depth}: G2\par \textit{Color}: none & \textit{K D}: (0,\allowbreak{}24,\allowbreak{}0)\allowbreak{}x1; (40,\allowbreak{}78,\allowbreak{}27)\allowbreak{}x2; (80,\allowbreak{}54,\allowbreak{}7)\allowbreak{}x1; (80,\allowbreak{}66,\allowbreak{}2)\allowbreak{}x1; (80,\allowbreak{}66,\allowbreak{}3)\allowbreak{}x1; (100,\allowbreak{}66,\allowbreak{}0)\allowbreak{}x1; (100,\allowbreak{}66,\allowbreak{}2)\allowbreak{}x1\par \textit{K Ω}: (80,\allowbreak{}54,\allowbreak{}6)\allowbreak{}x8 \\
% VI-CATALOG-ROW multisecciones 3
charged\_\allowbreak{}lepton\_\allowbreak{}G4 & \textit{Rank}: 16\par \textit{Cells (identifiers)}: 1,\allowbreak{}3,\allowbreak{}5,\allowbreak{}7,\allowbreak{}9,\allowbreak{}19,\allowbreak{}27,\allowbreak{}37,\allowbreak{}45,\allowbreak{}55,\allowbreak{}63,\allowbreak{}73,\allowbreak{}75,\allowbreak{}77,\allowbreak{}79,\allowbreak{}81\par \textit{Collections}: 16 & \textit{Collection}: charged\_\allowbreak{}lepton\par \textit{Depth}: G4\par \textit{Color}: none & \textit{K D}: (0,\allowbreak{}24,\allowbreak{}0)\allowbreak{}x1; (40,\allowbreak{}24,\allowbreak{}2)\allowbreak{}x2; (40,\allowbreak{}54,\allowbreak{}6)\allowbreak{}x2; (40,\allowbreak{}84,\allowbreak{}30)\allowbreak{}x2; (60,\allowbreak{}78,\allowbreak{}25)\allowbreak{}x3; (80,\allowbreak{}66,\allowbreak{}2)\allowbreak{}x2; (80,\allowbreak{}66,\allowbreak{}3)\allowbreak{}x3; (80,\allowbreak{}66,\allowbreak{}4)\allowbreak{}x1\par \textit{K Ω}: (24,\allowbreak{}54,\allowbreak{}2)\allowbreak{}x1; (56,\allowbreak{}54,\allowbreak{}5)\allowbreak{}x4; (80,\allowbreak{}54,\allowbreak{}6)\allowbreak{}x11 \\
% VI-CATALOG-ROW multisecciones 4
neutrino\_\allowbreak{}G1 & \textit{Rank}: 2\par \textit{Cells (identifiers)}: 40,\allowbreak{}42\par \textit{Collections}: 2 & \textit{Collection}: neutrino\par \textit{Depth}: G1\par \textit{Color}: none & \textit{K D}: (40,\allowbreak{}24,\allowbreak{}2)\allowbreak{}x1; (40,\allowbreak{}78,\allowbreak{}27)\allowbreak{}x1\par \textit{K Ω}: (40,\allowbreak{}54,\allowbreak{}4)\allowbreak{}x1; (80,\allowbreak{}54,\allowbreak{}6)\allowbreak{}x1 \\
% VI-CATALOG-ROW multisecciones 5
neutrino\_\allowbreak{}G2 & \textit{Rank}: 4\par \textit{Cells (identifiers)}: 22,\allowbreak{}24,\allowbreak{}58,\allowbreak{}60\par \textit{Collections}: 4 & \textit{Collection}: neutrino\par \textit{Depth}: G2\par \textit{Color}: none & \textit{K D}: (0,\allowbreak{}24,\allowbreak{}0)\allowbreak{}x1; (40,\allowbreak{}84,\allowbreak{}30)\allowbreak{}x1; (60,\allowbreak{}78,\allowbreak{}25)\allowbreak{}x1; (80,\allowbreak{}66,\allowbreak{}3)\allowbreak{}x1\par \textit{K Ω}: (4,\allowbreak{}54,\allowbreak{}2)\allowbreak{}x1; (80,\allowbreak{}54,\allowbreak{}6)\allowbreak{}x3 \\
% VI-CATALOG-ROW multisecciones 6
neutrino\_\allowbreak{}G3 & \textit{Rank}: 6\par \textit{Cells (identifiers)}: 20,\allowbreak{}26,\allowbreak{}38,\allowbreak{}44,\allowbreak{}56,\allowbreak{}62\par \textit{Collections}: 6 & \textit{Collection}: neutrino\par \textit{Depth}: G3\par \textit{Color}: none & \textit{K D}: (40,\allowbreak{}24,\allowbreak{}2)\allowbreak{}x2; (40,\allowbreak{}78,\allowbreak{}27)\allowbreak{}x1; (60,\allowbreak{}84,\allowbreak{}28)\allowbreak{}x1; (80,\allowbreak{}54,\allowbreak{}6)\allowbreak{}x1; (80,\allowbreak{}66,\allowbreak{}3)\allowbreak{}x1\par \textit{K Ω}: (36,\allowbreak{}54,\allowbreak{}4)\allowbreak{}x1; (56,\allowbreak{}54,\allowbreak{}5)\allowbreak{}x1; (80,\allowbreak{}54,\allowbreak{}6)\allowbreak{}x4 \\
% VI-CATALOG-ROW multisecciones 7
neutrino\_\allowbreak{}G4 & \textit{Rank}: 8\par \textit{Cells (identifiers)}: 2,\allowbreak{}4,\allowbreak{}6,\allowbreak{}8,\allowbreak{}74,\allowbreak{}76,\allowbreak{}78,\allowbreak{}80\par \textit{Collections}: 8 & \textit{Collection}: neutrino\par \textit{Depth}: G4\par \textit{Color}: none & \textit{K D}: (0,\allowbreak{}24,\allowbreak{}0)\allowbreak{}x1; (40,\allowbreak{}24,\allowbreak{}2)\allowbreak{}x1; (40,\allowbreak{}78,\allowbreak{}27)\allowbreak{}x2; (40,\allowbreak{}84,\allowbreak{}30)\allowbreak{}x1; (80,\allowbreak{}54,\allowbreak{}7)\allowbreak{}x1; (80,\allowbreak{}66,\allowbreak{}2)\allowbreak{}x1; (100,\allowbreak{}66,\allowbreak{}0)\allowbreak{}x1\par \textit{K Ω}: (36,\allowbreak{}54,\allowbreak{}4)\allowbreak{}x1; (56,\allowbreak{}54,\allowbreak{}5)\allowbreak{}x1; (80,\allowbreak{}54,\allowbreak{}6)\allowbreak{}x6 \\
% VI-CATALOG-ROW multisecciones 8
quark\_\allowbreak{}down\_\allowbreak{}G1 & \textit{Rank}: 2\par \textit{Cells (identifiers)}: 32,\allowbreak{}50\par \textit{Collections}: 2 & \textit{Collection}: quark\_\allowbreak{}down\par \textit{Depth}: G1\par \textit{Color}: B | G & \textit{K D}: (40,\allowbreak{}24,\allowbreak{}2)\allowbreak{}x1; (40,\allowbreak{}78,\allowbreak{}27)\allowbreak{}x1\par \textit{K Ω}: (40,\allowbreak{}54,\allowbreak{}4)\allowbreak{}x1; (80,\allowbreak{}54,\allowbreak{}6)\allowbreak{}x1 \\
% VI-CATALOG-ROW multisecciones 9
quark\_\allowbreak{}down\_\allowbreak{}G2 & \textit{Rank}: 4\par \textit{Cells (identifiers)}: 30,\allowbreak{}34,\allowbreak{}48,\allowbreak{}52\par \textit{Collections}: 4 & \textit{Collection}: quark\_\allowbreak{}down\par \textit{Depth}: G2\par \textit{Color}: B | G | R & \textit{K D}: (0,\allowbreak{}24,\allowbreak{}0)\allowbreak{}x1; (40,\allowbreak{}84,\allowbreak{}30)\allowbreak{}x1; (60,\allowbreak{}78,\allowbreak{}25)\allowbreak{}x1; (80,\allowbreak{}66,\allowbreak{}2)\allowbreak{}x1\par \textit{K Ω}: (4,\allowbreak{}54,\allowbreak{}2)\allowbreak{}x1; (80,\allowbreak{}54,\allowbreak{}6)\allowbreak{}x3 \\
% VI-CATALOG-ROW multisecciones 10
quark\_\allowbreak{}down\_\allowbreak{}G3 & \textit{Rank}: 6\par \textit{Cells (identifiers)}: 12,\allowbreak{}14,\allowbreak{}16,\allowbreak{}66,\allowbreak{}68,\allowbreak{}70\par \textit{Collections}: 6 & \textit{Collection}: quark\_\allowbreak{}down\par \textit{Depth}: G3\par \textit{Color}: B | G | R & \textit{K D}: (40,\allowbreak{}24,\allowbreak{}2)\allowbreak{}x2; (40,\allowbreak{}78,\allowbreak{}27)\allowbreak{}x1; (60,\allowbreak{}84,\allowbreak{}28)\allowbreak{}x1; (80,\allowbreak{}54,\allowbreak{}6)\allowbreak{}x1; (80,\allowbreak{}66,\allowbreak{}4)\allowbreak{}x1\par \textit{K Ω}: (36,\allowbreak{}54,\allowbreak{}4)\allowbreak{}x1; (80,\allowbreak{}54,\allowbreak{}6)\allowbreak{}x5 \\
% VI-CATALOG-ROW multisecciones 11
quark\_\allowbreak{}down\_\allowbreak{}G4 & \textit{Rank}: 8\par \textit{Cells (identifiers)}: 10,\allowbreak{}18,\allowbreak{}28,\allowbreak{}36,\allowbreak{}46,\allowbreak{}54,\allowbreak{}64,\allowbreak{}72\par \textit{Collections}: 8 & \textit{Collection}: quark\_\allowbreak{}down\par \textit{Depth}: G4\par \textit{Color}: B | G | R & \textit{K D}: (0,\allowbreak{}24,\allowbreak{}0)\allowbreak{}x1; (40,\allowbreak{}24,\allowbreak{}2)\allowbreak{}x1; (40,\allowbreak{}78,\allowbreak{}27)\allowbreak{}x2; (40,\allowbreak{}84,\allowbreak{}30)\allowbreak{}x1; (80,\allowbreak{}54,\allowbreak{}7)\allowbreak{}x1; (80,\allowbreak{}66,\allowbreak{}3)\allowbreak{}x1; (100,\allowbreak{}66,\allowbreak{}2)\allowbreak{}x1\par \textit{K Ω}: (36,\allowbreak{}54,\allowbreak{}4)\allowbreak{}x1; (80,\allowbreak{}54,\allowbreak{}6)\allowbreak{}x7 \\
% VI-CATALOG-ROW multisecciones 12
quark\_\allowbreak{}up\_\allowbreak{}G1 & \textit{Rank}: 4\par \textit{Cells (identifiers)}: 31,\allowbreak{}33,\allowbreak{}49,\allowbreak{}51\par \textit{Collections}: 3 & \textit{Collection}: quark\_\allowbreak{}up\par \textit{Depth}: G1\par \textit{Color}: B | G | R & \textit{K D}: (40,\allowbreak{}54,\allowbreak{}6)\allowbreak{}x1; (60,\allowbreak{}78,\allowbreak{}25)\allowbreak{}x1; (80,\allowbreak{}66,\allowbreak{}2)\allowbreak{}x1; (80,\allowbreak{}66,\allowbreak{}3)\allowbreak{}x1\par \textit{K Ω}: (80,\allowbreak{}54,\allowbreak{}6)\allowbreak{}x4 \\
% VI-CATALOG-ROW multisecciones 13
quark\_\allowbreak{}up\_\allowbreak{}G3 & \textit{Rank}: 12\par \textit{Cells (identifiers)}: 11,\allowbreak{}13,\allowbreak{}15,\allowbreak{}17,\allowbreak{}29,\allowbreak{}35,\allowbreak{}47,\allowbreak{}53,\allowbreak{}65,\allowbreak{}67,\allowbreak{}69,\allowbreak{}71\par \textit{Collections}: 12 & \textit{Collection}: quark\_\allowbreak{}up\par \textit{Depth}: G3\par \textit{Color}: B | G | R & \textit{K D}: (40,\allowbreak{}24,\allowbreak{}2)\allowbreak{}x2; (40,\allowbreak{}78,\allowbreak{}27)\allowbreak{}x2; (60,\allowbreak{}84,\allowbreak{}28)\allowbreak{}x1; (80,\allowbreak{}54,\allowbreak{}6)\allowbreak{}x3; (80,\allowbreak{}66,\allowbreak{}3)\allowbreak{}x1; (80,\allowbreak{}66,\allowbreak{}4)\allowbreak{}x1; (100,\allowbreak{}66,\allowbreak{}0)\allowbreak{}x1; (100,\allowbreak{}66,\allowbreak{}2)\allowbreak{}x1\par \textit{K Ω}: (56,\allowbreak{}54,\allowbreak{}5)\allowbreak{}x3; (80,\allowbreak{}54,\allowbreak{}6)\allowbreak{}x9 \\
\end{longtable}
\endgroup
\subsection{The 56 route configurations}

The signature identifying a collection is preserved literally, with its signs
and separators. The cells and multisections make explicit that a route
collection is not, by definition, equivalent to a physical species. The joint
profiles count the combinations recorded within the collection. The decimal
spectra of the character and the complete signature spectra remain intact
in \texttt{catalogo\_familias.json}; the incidence and integer readers
expressed here provide a means of navigating that information.

\begingroup
\setlength{\tabcolsep}{4pt}
\setlength{\extrarowheight}{4pt}
\renewcommand{\arraystretch}{1.13}
\begin{longtable}{@{}>{\raggedright\arraybackslash}p{\dimexpr 0.21794872\linewidth-2\tabcolsep\relax}>{\raggedright\arraybackslash}p{\dimexpr 0.34615385\linewidth-2\tabcolsep\relax}>{\raggedright\arraybackslash}p{\dimexpr 0.21794872\linewidth-2\tabcolsep\relax}>{\raggedright\arraybackslash}p{\dimexpr 0.21794872\linewidth-2\tabcolsep\relax}@{}}
\caption{The 56 internal route configurations.}\label{tab:vi-familias}\\
\toprule
Collection and signature & Incidence & \(K_D\) & \(K_\Omega\) \\
\midrule
\endfirsthead
\multicolumn{4}{l}{\textit{Continuation}}\\
\toprule
Collection and signature & Incidence & \(K_D\) & \(K_\Omega\) \\
\midrule
\endhead
\midrule
\multicolumn{4}{r}{Continued on the next page}\\
\endfoot
\bottomrule
\endlastfoot
% VI-CATALOG-ROW familias 1
1\par 24|\allowbreak{}0|\allowbreak{}12|\allowbreak{}0|\allowbreak{}-12|\allowbreak{}+++-\kern0pt{}-\kern0pt{}-+++-\kern0pt{}-\kern0pt{}- & \textit{Number of cells}: 5\par \textit{Cells (identifiers)}: 21,\allowbreak{}31,\allowbreak{}41,\allowbreak{}51,\allowbreak{}81\par \textit{Multisections}: charged\_\allowbreak{}lepton\_\allowbreak{}G0 | charged\_\allowbreak{}lepton\_\allowbreak{}G2 | charged\_\allowbreak{}lepton\_\allowbreak{}G4 | quark\_\allowbreak{}up\_\allowbreak{}G1\par \textit{Collections}: charged\_\allowbreak{}lepton | quark\_\allowbreak{}up\par \textit{Color}: R | none\par \textit{Joint profiles}: 5 & (0,\allowbreak{}24,\allowbreak{}0)\allowbreak{}x2; (40,\allowbreak{}54,\allowbreak{}6)\allowbreak{}x1; (60,\allowbreak{}78,\allowbreak{}25)\allowbreak{}x1; (80,\allowbreak{}54,\allowbreak{}6)\allowbreak{}x1 & (24,\allowbreak{}54,\allowbreak{}2)\allowbreak{}x1; (80,\allowbreak{}54,\allowbreak{}6)\allowbreak{}x4 \\
% VI-CATALOG-ROW familias 2
2\par 24|\allowbreak{}0|\allowbreak{}4|\allowbreak{}0|\allowbreak{}-12|\allowbreak{}+++-\kern0pt{}-\kern0pt{}-+++-\kern0pt{}-\kern0pt{}- & \textit{Number of cells}: 2\par \textit{Cells (identifiers)}: 11,\allowbreak{}61\par \textit{Multisections}: charged\_\allowbreak{}lepton\_\allowbreak{}G2 | quark\_\allowbreak{}up\_\allowbreak{}G3\par \textit{Collections}: charged\_\allowbreak{}lepton | quark\_\allowbreak{}up\par \textit{Color}: R | none\par \textit{Joint profiles}: 2 & (60,\allowbreak{}84,\allowbreak{}28)\allowbreak{}x1; (80,\allowbreak{}54,\allowbreak{}7)\allowbreak{}x1 & (80,\allowbreak{}54,\allowbreak{}6)\allowbreak{}x2 \\
% VI-CATALOG-ROW familias 3
3\par 24|\allowbreak{}0|\allowbreak{}8|\allowbreak{}0|\allowbreak{}-12|\allowbreak{}+++-\kern0pt{}-\kern0pt{}-+++-\kern0pt{}-\kern0pt{}- & \textit{Number of cells}: 2\par \textit{Cells (identifiers)}: 1,\allowbreak{}71\par \textit{Multisections}: charged\_\allowbreak{}lepton\_\allowbreak{}G4 | quark\_\allowbreak{}up\_\allowbreak{}G3\par \textit{Collections}: charged\_\allowbreak{}lepton | quark\_\allowbreak{}up\par \textit{Color}: R | none\par \textit{Joint profiles}: 2 & (60,\allowbreak{}78,\allowbreak{}25)\allowbreak{}x1; (80,\allowbreak{}54,\allowbreak{}6)\allowbreak{}x1 & (80,\allowbreak{}54,\allowbreak{}6)\allowbreak{}x2 \\
% VI-CATALOG-ROW familias 4
4\par 28|\allowbreak{}-6|\allowbreak{}-2|\allowbreak{}-2|\allowbreak{}-8|\allowbreak{}0++-\kern0pt{}-\kern0pt{}-0++-\kern0pt{}-\kern0pt{}- & \textit{Number of cells}: 1\par \textit{Cells (identifiers)}: 53\par \textit{Multisections}: quark\_\allowbreak{}up\_\allowbreak{}G3\par \textit{Collections}: quark\_\allowbreak{}up\par \textit{Color}: G\par \textit{Joint profiles}: 1 & (80,\allowbreak{}66,\allowbreak{}3)\allowbreak{}x1 & (56,\allowbreak{}54,\allowbreak{}5)\allowbreak{}x1 \\
% VI-CATALOG-ROW familias 5
5\par 28|\allowbreak{}-6|\allowbreak{}-4|\allowbreak{}-2|\allowbreak{}-8|\allowbreak{}0++-\kern0pt{}-\kern0pt{}-0++-\kern0pt{}-\kern0pt{}- & \textit{Number of cells}: 1\par \textit{Cells (identifiers)}: 20\par \textit{Multisections}: neutrino\_\allowbreak{}G3\par \textit{Collections}: neutrino\par \textit{Color}: none\par \textit{Joint profiles}: 1 & (80,\allowbreak{}66,\allowbreak{}3)\allowbreak{}x1 & (56,\allowbreak{}54,\allowbreak{}5)\allowbreak{}x1 \\
% VI-CATALOG-ROW familias 6
6\par 28|\allowbreak{}-6|\allowbreak{}0|\allowbreak{}0|\allowbreak{}-12|\allowbreak{}+++-\kern0pt{}-\kern0pt{}-+++-\kern0pt{}-\kern0pt{}- & \textit{Number of cells}: 2\par \textit{Cells (identifiers)}: 30,\allowbreak{}70\par \textit{Multisections}: quark\_\allowbreak{}down\_\allowbreak{}G2 | quark\_\allowbreak{}down\_\allowbreak{}G3\par \textit{Collections}: quark\_\allowbreak{}down\par \textit{Color}: G\par \textit{Joint profiles}: 2 & (40,\allowbreak{}24,\allowbreak{}2)\allowbreak{}x1; (80,\allowbreak{}66,\allowbreak{}2)\allowbreak{}x1 & (80,\allowbreak{}54,\allowbreak{}6)\allowbreak{}x2 \\
% VI-CATALOG-ROW familias 7
7\par 28|\allowbreak{}-6|\allowbreak{}10|\allowbreak{}0|\allowbreak{}-12|\allowbreak{}+++-\kern0pt{}-\kern0pt{}-+++-\kern0pt{}-\kern0pt{}- & \textit{Number of cells}: 2\par \textit{Cells (identifiers)}: 33,\allowbreak{}64\par \textit{Multisections}: quark\_\allowbreak{}down\_\allowbreak{}G4 | quark\_\allowbreak{}up\_\allowbreak{}G1\par \textit{Collections}: quark\_\allowbreak{}down | quark\_\allowbreak{}up\par \textit{Color}: G\par \textit{Joint profiles}: 2 & (40,\allowbreak{}24,\allowbreak{}2)\allowbreak{}x1; (80,\allowbreak{}66,\allowbreak{}2)\allowbreak{}x1 & (36,\allowbreak{}54,\allowbreak{}4)\allowbreak{}x1; (80,\allowbreak{}54,\allowbreak{}6)\allowbreak{}x1 \\
% VI-CATALOG-ROW familias 8
8\par 28|\allowbreak{}-6|\allowbreak{}12|\allowbreak{}0|\allowbreak{}-12|\allowbreak{}+++-\kern0pt{}-\kern0pt{}-+++-\kern0pt{}-\kern0pt{}- & \textit{Number of cells}: 1\par \textit{Cells (identifiers)}: 10\par \textit{Multisections}: quark\_\allowbreak{}down\_\allowbreak{}G4\par \textit{Collections}: quark\_\allowbreak{}down\par \textit{Color}: G\par \textit{Joint profiles}: 1 & (100,\allowbreak{}66,\allowbreak{}2)\allowbreak{}x1 & (80,\allowbreak{}54,\allowbreak{}6)\allowbreak{}x1 \\
% VI-CATALOG-ROW familias 9
9\par 28|\allowbreak{}-6|\allowbreak{}2|\allowbreak{}0|\allowbreak{}-12|\allowbreak{}+++-\kern0pt{}-\kern0pt{}-+++-\kern0pt{}-\kern0pt{}- & \textit{Number of cells}: 1\par \textit{Cells (identifiers)}: 3\par \textit{Multisections}: charged\_\allowbreak{}lepton\_\allowbreak{}G4\par \textit{Collections}: charged\_\allowbreak{}lepton\par \textit{Color}: none\par \textit{Joint profiles}: 1 & (40,\allowbreak{}84,\allowbreak{}30)\allowbreak{}x1 & (80,\allowbreak{}54,\allowbreak{}6)\allowbreak{}x1 \\
% VI-CATALOG-ROW familias 10
10\par 28|\allowbreak{}-6|\allowbreak{}4|\allowbreak{}-2|\allowbreak{}-8|\allowbreak{}++0-\kern0pt{}-\kern0pt{}-++0-\kern0pt{}-\kern0pt{}- & \textit{Number of cells}: 1\par \textit{Cells (identifiers)}: 50\par \textit{Multisections}: quark\_\allowbreak{}down\_\allowbreak{}G1\par \textit{Collections}: quark\_\allowbreak{}down\par \textit{Color}: G\par \textit{Joint profiles}: 1 & (40,\allowbreak{}78,\allowbreak{}27)\allowbreak{}x1 & (80,\allowbreak{}54,\allowbreak{}6)\allowbreak{}x1 \\
% VI-CATALOG-ROW familias 11
11\par 28|\allowbreak{}-6|\allowbreak{}4|\allowbreak{}0|\allowbreak{}-12|\allowbreak{}+++-\kern0pt{}-\kern0pt{}-+++-\kern0pt{}-\kern0pt{}- & \textit{Number of cells}: 2\par \textit{Cells (identifiers)}: 9,\allowbreak{}40\par \textit{Multisections}: charged\_\allowbreak{}lepton\_\allowbreak{}G4 | neutrino\_\allowbreak{}G1\par \textit{Collections}: charged\_\allowbreak{}lepton | neutrino\par \textit{Color}: none\par \textit{Joint profiles}: 2 & (40,\allowbreak{}24,\allowbreak{}2)\allowbreak{}x1; (80,\allowbreak{}66,\allowbreak{}2)\allowbreak{}x1 & (40,\allowbreak{}54,\allowbreak{}4)\allowbreak{}x1; (80,\allowbreak{}54,\allowbreak{}6)\allowbreak{}x1 \\
% VI-CATALOG-ROW familias 12
12\par 28|\allowbreak{}-6|\allowbreak{}4|\allowbreak{}0|\allowbreak{}-8|\allowbreak{}+0+-\kern0pt{}-\kern0pt{}-+0+-\kern0pt{}-\kern0pt{}- & \textit{Number of cells}: 1\par \textit{Cells (identifiers)}: 80\par \textit{Multisections}: neutrino\_\allowbreak{}G4\par \textit{Collections}: neutrino\par \textit{Color}: none\par \textit{Joint profiles}: 1 & (40,\allowbreak{}78,\allowbreak{}27)\allowbreak{}x1 & (80,\allowbreak{}54,\allowbreak{}6)\allowbreak{}x1 \\
% VI-CATALOG-ROW familias 13
13\par 28|\allowbreak{}-6|\allowbreak{}6|\allowbreak{}-2|\allowbreak{}-8|\allowbreak{}++0-\kern0pt{}-\kern0pt{}-++0-\kern0pt{}-\kern0pt{}- & \textit{Number of cells}: 1\par \textit{Cells (identifiers)}: 74\par \textit{Multisections}: neutrino\_\allowbreak{}G4\par \textit{Collections}: neutrino\par \textit{Color}: none\par \textit{Joint profiles}: 1 & (40,\allowbreak{}78,\allowbreak{}27)\allowbreak{}x1 & (80,\allowbreak{}54,\allowbreak{}6)\allowbreak{}x1 \\
% VI-CATALOG-ROW familias 14
14\par 28|\allowbreak{}-6|\allowbreak{}6|\allowbreak{}0|\allowbreak{}-12|\allowbreak{}+++-\kern0pt{}-\kern0pt{}-+++-\kern0pt{}-\kern0pt{}- & \textit{Number of cells}: 3\par \textit{Cells (identifiers)}: 13,\allowbreak{}43,\allowbreak{}63\par \textit{Multisections}: charged\_\allowbreak{}lepton\_\allowbreak{}G2 | charged\_\allowbreak{}lepton\_\allowbreak{}G4 | quark\_\allowbreak{}up\_\allowbreak{}G3\par \textit{Collections}: charged\_\allowbreak{}lepton | quark\_\allowbreak{}up\par \textit{Color}: G | none\par \textit{Joint profiles}: 3 & (40,\allowbreak{}24,\allowbreak{}2)\allowbreak{}x1; (80,\allowbreak{}66,\allowbreak{}2)\allowbreak{}x1; (100,\allowbreak{}66,\allowbreak{}2)\allowbreak{}x1 & (80,\allowbreak{}54,\allowbreak{}6)\allowbreak{}x3 \\
% VI-CATALOG-ROW familias 15
15\par 28|\allowbreak{}-6|\allowbreak{}6|\allowbreak{}0|\allowbreak{}-8|\allowbreak{}+0+-\kern0pt{}-\kern0pt{}-+0+-\kern0pt{}-\kern0pt{}- & \textit{Number of cells}: 1\par \textit{Cells (identifiers)}: 23\par \textit{Multisections}: charged\_\allowbreak{}lepton\_\allowbreak{}G2\par \textit{Collections}: charged\_\allowbreak{}lepton\par \textit{Color}: none\par \textit{Joint profiles}: 1 & (40,\allowbreak{}78,\allowbreak{}27)\allowbreak{}x1 & (80,\allowbreak{}54,\allowbreak{}6)\allowbreak{}x1 \\
% VI-CATALOG-ROW familias 16
16\par 28|\allowbreak{}-6|\allowbreak{}8|\allowbreak{}0|\allowbreak{}-12|\allowbreak{}+++-\kern0pt{}-\kern0pt{}-+++-\kern0pt{}-\kern0pt{}- & \textit{Number of cells}: 1\par \textit{Cells (identifiers)}: 60\par \textit{Multisections}: neutrino\_\allowbreak{}G2\par \textit{Collections}: neutrino\par \textit{Color}: none\par \textit{Joint profiles}: 1 & (40,\allowbreak{}84,\allowbreak{}30)\allowbreak{}x1 & (80,\allowbreak{}54,\allowbreak{}6)\allowbreak{}x1 \\
% VI-CATALOG-ROW familias 17
17\par 28|\allowbreak{}6|\allowbreak{}-10|\allowbreak{}0|\allowbreak{}-12|\allowbreak{}+++-\kern0pt{}-\kern0pt{}-+++-\kern0pt{}-\kern0pt{}- & \textit{Number of cells}: 1\par \textit{Cells (identifiers)}: 19\par \textit{Multisections}: charged\_\allowbreak{}lepton\_\allowbreak{}G4\par \textit{Collections}: charged\_\allowbreak{}lepton\par \textit{Color}: none\par \textit{Joint profiles}: 1 & (40,\allowbreak{}84,\allowbreak{}30)\allowbreak{}x1 & (80,\allowbreak{}54,\allowbreak{}6)\allowbreak{}x1 \\
% VI-CATALOG-ROW familias 18
18\par 28|\allowbreak{}6|\allowbreak{}-2|\allowbreak{}0|\allowbreak{}-12|\allowbreak{}+++-\kern0pt{}-\kern0pt{}-+++-\kern0pt{}-\kern0pt{}- & \textit{Number of cells}: 5\par \textit{Cells (identifiers)}: 18,\allowbreak{}29,\allowbreak{}49,\allowbreak{}59,\allowbreak{}79\par \textit{Multisections}: charged\_\allowbreak{}lepton\_\allowbreak{}G2 | charged\_\allowbreak{}lepton\_\allowbreak{}G4 | quark\_\allowbreak{}down\_\allowbreak{}G4 | quark\_\allowbreak{}up\_\allowbreak{}G1 | quark\_\allowbreak{}up\_\allowbreak{}G3\par \textit{Collections}: charged\_\allowbreak{}lepton | quark\_\allowbreak{}down | quark\_\allowbreak{}up\par \textit{Color}: B | none\par \textit{Joint profiles}: 4 & (40,\allowbreak{}24,\allowbreak{}2)\allowbreak{}x1; (40,\allowbreak{}78,\allowbreak{}27)\allowbreak{}x1; (80,\allowbreak{}66,\allowbreak{}3)\allowbreak{}x2; (100,\allowbreak{}66,\allowbreak{}0)\allowbreak{}x1 & (80,\allowbreak{}54,\allowbreak{}6)\allowbreak{}x5 \\
% VI-CATALOG-ROW familias 19
19\par 28|\allowbreak{}6|\allowbreak{}-4|\allowbreak{}0|\allowbreak{}-12|\allowbreak{}+++-\kern0pt{}-\kern0pt{}-+++-\kern0pt{}-\kern0pt{}- & \textit{Number of cells}: 2\par \textit{Cells (identifiers)}: 2,\allowbreak{}42\par \textit{Multisections}: neutrino\_\allowbreak{}G1 | neutrino\_\allowbreak{}G4\par \textit{Collections}: neutrino\par \textit{Color}: none\par \textit{Joint profiles}: 2 & (40,\allowbreak{}78,\allowbreak{}27)\allowbreak{}x1; (100,\allowbreak{}66,\allowbreak{}0)\allowbreak{}x1 & (80,\allowbreak{}54,\allowbreak{}6)\allowbreak{}x2 \\
% VI-CATALOG-ROW familias 20
20\par 28|\allowbreak{}6|\allowbreak{}-6|\allowbreak{}0|\allowbreak{}-12|\allowbreak{}+++-\kern0pt{}-\kern0pt{}-+++-\kern0pt{}-\kern0pt{}- & \textit{Number of cells}: 1\par \textit{Cells (identifiers)}: 69\par \textit{Multisections}: quark\_\allowbreak{}up\_\allowbreak{}G3\par \textit{Collections}: quark\_\allowbreak{}up\par \textit{Color}: B\par \textit{Joint profiles}: 1 & (80,\allowbreak{}66,\allowbreak{}4)\allowbreak{}x1 & (80,\allowbreak{}54,\allowbreak{}6)\allowbreak{}x1 \\
% VI-CATALOG-ROW familias 21
21\par 28|\allowbreak{}6|\allowbreak{}-8|\allowbreak{}0|\allowbreak{}-12|\allowbreak{}+++-\kern0pt{}-\kern0pt{}-+++-\kern0pt{}-\kern0pt{}- & \textit{Number of cells}: 1\par \textit{Cells (identifiers)}: 52\par \textit{Multisections}: quark\_\allowbreak{}down\_\allowbreak{}G2\par \textit{Collections}: quark\_\allowbreak{}down\par \textit{Color}: B\par \textit{Joint profiles}: 1 & (40,\allowbreak{}84,\allowbreak{}30)\allowbreak{}x1 & (80,\allowbreak{}54,\allowbreak{}6)\allowbreak{}x1 \\
% VI-CATALOG-ROW familias 22
22\par 28|\allowbreak{}6|\allowbreak{}0|\allowbreak{}0|\allowbreak{}-12|\allowbreak{}+++-\kern0pt{}-\kern0pt{}-+++-\kern0pt{}-\kern0pt{}- & \textit{Number of cells}: 4\par \textit{Cells (identifiers)}: 22,\allowbreak{}32,\allowbreak{}72,\allowbreak{}73\par \textit{Multisections}: charged\_\allowbreak{}lepton\_\allowbreak{}G4 | neutrino\_\allowbreak{}G2 | quark\_\allowbreak{}down\_\allowbreak{}G1 | quark\_\allowbreak{}down\_\allowbreak{}G4\par \textit{Collections}: charged\_\allowbreak{}lepton | neutrino | quark\_\allowbreak{}down\par \textit{Color}: B | none\par \textit{Joint profiles}: 3 & (40,\allowbreak{}24,\allowbreak{}2)\allowbreak{}x1; (40,\allowbreak{}78,\allowbreak{}27)\allowbreak{}x1; (80,\allowbreak{}66,\allowbreak{}3)\allowbreak{}x2 & (40,\allowbreak{}54,\allowbreak{}4)\allowbreak{}x1; (80,\allowbreak{}54,\allowbreak{}6)\allowbreak{}x3 \\
% VI-CATALOG-ROW familias 23
23\par 28|\allowbreak{}6|\allowbreak{}2|\allowbreak{}0|\allowbreak{}-12|\allowbreak{}+++-\kern0pt{}-\kern0pt{}-+++-\kern0pt{}-\kern0pt{}- & \textit{Number of cells}: 2\par \textit{Cells (identifiers)}: 8,\allowbreak{}39\par \textit{Multisections}: charged\_\allowbreak{}lepton\_\allowbreak{}G2 | neutrino\_\allowbreak{}G4\par \textit{Collections}: charged\_\allowbreak{}lepton | neutrino\par \textit{Color}: none\par \textit{Joint profiles}: 2 & (40,\allowbreak{}24,\allowbreak{}2)\allowbreak{}x1; (40,\allowbreak{}78,\allowbreak{}27)\allowbreak{}x1 & (36,\allowbreak{}54,\allowbreak{}4)\allowbreak{}x1; (80,\allowbreak{}54,\allowbreak{}6)\allowbreak{}x1 \\
% VI-CATALOG-ROW familias 24
24\par 28|\allowbreak{}6|\allowbreak{}4|\allowbreak{}0|\allowbreak{}-12|\allowbreak{}+++-\kern0pt{}-\kern0pt{}-+++-\kern0pt{}-\kern0pt{}- & \textit{Number of cells}: 1\par \textit{Cells (identifiers)}: 12\par \textit{Multisections}: quark\_\allowbreak{}down\_\allowbreak{}G3\par \textit{Collections}: quark\_\allowbreak{}down\par \textit{Color}: B\par \textit{Joint profiles}: 1 & (80,\allowbreak{}66,\allowbreak{}4)\allowbreak{}x1 & (80,\allowbreak{}54,\allowbreak{}6)\allowbreak{}x1 \\
% VI-CATALOG-ROW familias 25
25\par 28|\allowbreak{}6|\allowbreak{}8|\allowbreak{}0|\allowbreak{}-12|\allowbreak{}+++-\kern0pt{}-\kern0pt{}-+++-\kern0pt{}-\kern0pt{}- & \textit{Number of cells}: 1\par \textit{Cells (identifiers)}: 62\par \textit{Multisections}: neutrino\_\allowbreak{}G3\par \textit{Collections}: neutrino\par \textit{Color}: none\par \textit{Joint profiles}: 1 & (40,\allowbreak{}24,\allowbreak{}2)\allowbreak{}x1 & (80,\allowbreak{}54,\allowbreak{}6)\allowbreak{}x1 \\
% VI-CATALOG-ROW familias 26
26\par 30|\allowbreak{}-6|\allowbreak{}-4|\allowbreak{}-2|\allowbreak{}-8|\allowbreak{}0++-\kern0pt{}-\kern0pt{}-0++-\kern0pt{}-\kern0pt{}- & \textit{Number of cells}: 1\par \textit{Cells (identifiers)}: 4\par \textit{Multisections}: neutrino\_\allowbreak{}G4\par \textit{Collections}: neutrino\par \textit{Color}: none\par \textit{Joint profiles}: 1 & (80,\allowbreak{}54,\allowbreak{}7)\allowbreak{}x1 & (56,\allowbreak{}54,\allowbreak{}5)\allowbreak{}x1 \\
% VI-CATALOG-ROW familias 27
27\par 30|\allowbreak{}-6|\allowbreak{}-8|\allowbreak{}0|\allowbreak{}-8|\allowbreak{}+0+-\kern0pt{}-\kern0pt{}-+0+-\kern0pt{}-\kern0pt{}- & \textit{Number of cells}: 1\par \textit{Cells (identifiers)}: 55\par \textit{Multisections}: charged\_\allowbreak{}lepton\_\allowbreak{}G4\par \textit{Collections}: charged\_\allowbreak{}lepton\par \textit{Color}: none\par \textit{Joint profiles}: 1 & (60,\allowbreak{}78,\allowbreak{}25)\allowbreak{}x1 & (80,\allowbreak{}54,\allowbreak{}6)\allowbreak{}x1 \\
% VI-CATALOG-ROW familias 28
28\par 30|\allowbreak{}-6|\allowbreak{}0|\allowbreak{}-2|\allowbreak{}-8|\allowbreak{}++0-\kern0pt{}-\kern0pt{}-++0-\kern0pt{}-\kern0pt{}- & \textit{Number of cells}: 2\par \textit{Cells (identifiers)}: 24,\allowbreak{}34\par \textit{Multisections}: neutrino\_\allowbreak{}G2 | quark\_\allowbreak{}down\_\allowbreak{}G2\par \textit{Collections}: neutrino | quark\_\allowbreak{}down\par \textit{Color}: R | none\par \textit{Joint profiles}: 2 & (0,\allowbreak{}24,\allowbreak{}0)\allowbreak{}x1; (60,\allowbreak{}78,\allowbreak{}25)\allowbreak{}x1 & (4,\allowbreak{}54,\allowbreak{}2)\allowbreak{}x1; (80,\allowbreak{}54,\allowbreak{}6)\allowbreak{}x1 \\
% VI-CATALOG-ROW familias 29
29\par 30|\allowbreak{}-6|\allowbreak{}0|\allowbreak{}-2|\allowbreak{}-8|\allowbreak{}0++-\kern0pt{}-\kern0pt{}-0++-\kern0pt{}-\kern0pt{}- & \textit{Number of cells}: 1\par \textit{Cells (identifiers)}: 75\par \textit{Multisections}: charged\_\allowbreak{}lepton\_\allowbreak{}G4\par \textit{Collections}: charged\_\allowbreak{}lepton\par \textit{Color}: none\par \textit{Joint profiles}: 1 & (40,\allowbreak{}54,\allowbreak{}6)\allowbreak{}x1 & (56,\allowbreak{}54,\allowbreak{}5)\allowbreak{}x1 \\
% VI-CATALOG-ROW familias 30
30\par 30|\allowbreak{}-6|\allowbreak{}0|\allowbreak{}0|\allowbreak{}-12|\allowbreak{}+++-\kern0pt{}-\kern0pt{}-+++-\kern0pt{}-\kern0pt{}- & \textit{Number of cells}: 1\par \textit{Cells (identifiers)}: 14\par \textit{Multisections}: quark\_\allowbreak{}down\_\allowbreak{}G3\par \textit{Collections}: quark\_\allowbreak{}down\par \textit{Color}: R\par \textit{Joint profiles}: 1 & (80,\allowbreak{}54,\allowbreak{}6)\allowbreak{}x1 & (80,\allowbreak{}54,\allowbreak{}6)\allowbreak{}x1 \\
% VI-CATALOG-ROW familias 31
31\par 30|\allowbreak{}-6|\allowbreak{}0|\allowbreak{}0|\allowbreak{}-8|\allowbreak{}+0+-\kern0pt{}-\kern0pt{}-+0+-\kern0pt{}-\kern0pt{}- & \textit{Number of cells}: 1\par \textit{Cells (identifiers)}: 54\par \textit{Multisections}: quark\_\allowbreak{}down\_\allowbreak{}G4\par \textit{Collections}: quark\_\allowbreak{}down\par \textit{Color}: R\par \textit{Joint profiles}: 1 & (0,\allowbreak{}24,\allowbreak{}0)\allowbreak{}x1 & (80,\allowbreak{}54,\allowbreak{}6)\allowbreak{}x1 \\
% VI-CATALOG-ROW familias 32
32\par 30|\allowbreak{}-6|\allowbreak{}4|\allowbreak{}0|\allowbreak{}-12|\allowbreak{}+++-\kern0pt{}-\kern0pt{}-+++-\kern0pt{}-\kern0pt{}- & \textit{Number of cells}: 1\par \textit{Cells (identifiers)}: 65\par \textit{Multisections}: quark\_\allowbreak{}up\_\allowbreak{}G3\par \textit{Collections}: quark\_\allowbreak{}up\par \textit{Color}: R\par \textit{Joint profiles}: 1 & (80,\allowbreak{}54,\allowbreak{}6)\allowbreak{}x1 & (80,\allowbreak{}54,\allowbreak{}6)\allowbreak{}x1 \\
% VI-CATALOG-ROW familias 33
33\par 30|\allowbreak{}-6|\allowbreak{}8|\allowbreak{}0|\allowbreak{}-12|\allowbreak{}+++-\kern0pt{}-\kern0pt{}-+++-\kern0pt{}-\kern0pt{}- & \textit{Number of cells}: 1\par \textit{Cells (identifiers)}: 44\par \textit{Multisections}: neutrino\_\allowbreak{}G3\par \textit{Collections}: neutrino\par \textit{Color}: none\par \textit{Joint profiles}: 1 & (60,\allowbreak{}84,\allowbreak{}28)\allowbreak{}x1 & (80,\allowbreak{}54,\allowbreak{}6)\allowbreak{}x1 \\
% VI-CATALOG-ROW familias 34
34\par 30|\allowbreak{}6|\allowbreak{}-4|\allowbreak{}0|\allowbreak{}-12|\allowbreak{}+++-\kern0pt{}-\kern0pt{}-+++-\kern0pt{}-\kern0pt{}- & \textit{Number of cells}: 1\par \textit{Cells (identifiers)}: 68\par \textit{Multisections}: quark\_\allowbreak{}down\_\allowbreak{}G3\par \textit{Collections}: quark\_\allowbreak{}down\par \textit{Color}: R\par \textit{Joint profiles}: 1 & (60,\allowbreak{}84,\allowbreak{}28)\allowbreak{}x1 & (80,\allowbreak{}54,\allowbreak{}6)\allowbreak{}x1 \\
% VI-CATALOG-ROW familias 35
35\par 30|\allowbreak{}6|\allowbreak{}-8|\allowbreak{}0|\allowbreak{}-12|\allowbreak{}+++-\kern0pt{}-\kern0pt{}-+++-\kern0pt{}-\kern0pt{}- & \textit{Number of cells}: 1\par \textit{Cells (identifiers)}: 17\par \textit{Multisections}: quark\_\allowbreak{}up\_\allowbreak{}G3\par \textit{Collections}: quark\_\allowbreak{}up\par \textit{Color}: R\par \textit{Joint profiles}: 1 & (80,\allowbreak{}54,\allowbreak{}6)\allowbreak{}x1 & (80,\allowbreak{}54,\allowbreak{}6)\allowbreak{}x1 \\
% VI-CATALOG-ROW familias 36
36\par 30|\allowbreak{}6|\allowbreak{}0|\allowbreak{}0|\allowbreak{}-12|\allowbreak{}+++-\kern0pt{}-\kern0pt{}-+++-\kern0pt{}-\kern0pt{}- & \textit{Number of cells}: 2\par \textit{Cells (identifiers)}: 38,\allowbreak{}58\par \textit{Multisections}: neutrino\_\allowbreak{}G2 | neutrino\_\allowbreak{}G3\par \textit{Collections}: neutrino\par \textit{Color}: none\par \textit{Joint profiles}: 2 & (60,\allowbreak{}78,\allowbreak{}25)\allowbreak{}x1; (80,\allowbreak{}54,\allowbreak{}6)\allowbreak{}x1 & (80,\allowbreak{}54,\allowbreak{}6)\allowbreak{}x2 \\
% VI-CATALOG-ROW familias 37
37\par 30|\allowbreak{}6|\allowbreak{}0|\allowbreak{}0|\allowbreak{}-8|\allowbreak{}+++-0-+++-0- & \textit{Number of cells}: 1\par \textit{Cells (identifiers)}: 78\par \textit{Multisections}: neutrino\_\allowbreak{}G4\par \textit{Collections}: neutrino\par \textit{Color}: none\par \textit{Joint profiles}: 1 & (0,\allowbreak{}24,\allowbreak{}0)\allowbreak{}x1 & (80,\allowbreak{}54,\allowbreak{}6)\allowbreak{}x1 \\
% VI-CATALOG-ROW familias 38
38\par 30|\allowbreak{}6|\allowbreak{}0|\allowbreak{}2|\allowbreak{}-8|\allowbreak{}+++-\kern0pt{}-0+++-\kern0pt{}-0 & \textit{Number of cells}: 1\par \textit{Cells (identifiers)}: 27\par \textit{Multisections}: charged\_\allowbreak{}lepton\_\allowbreak{}G4\par \textit{Collections}: charged\_\allowbreak{}lepton\par \textit{Color}: none\par \textit{Joint profiles}: 1 & (40,\allowbreak{}54,\allowbreak{}6)\allowbreak{}x1 & (56,\allowbreak{}54,\allowbreak{}5)\allowbreak{}x1 \\
% VI-CATALOG-ROW familias 39
39\par 30|\allowbreak{}6|\allowbreak{}0|\allowbreak{}2|\allowbreak{}-8|\allowbreak{}+++0-\kern0pt{}-+++0-\kern0pt{}- & \textit{Number of cells}: 1\par \textit{Cells (identifiers)}: 48\par \textit{Multisections}: quark\_\allowbreak{}down\_\allowbreak{}G2\par \textit{Collections}: quark\_\allowbreak{}down\par \textit{Color}: R\par \textit{Joint profiles}: 1 & (0,\allowbreak{}24,\allowbreak{}0)\allowbreak{}x1 & (4,\allowbreak{}54,\allowbreak{}2)\allowbreak{}x1 \\
% VI-CATALOG-ROW familias 40
40\par 30|\allowbreak{}6|\allowbreak{}4|\allowbreak{}0|\allowbreak{}-12|\allowbreak{}+++-\kern0pt{}-\kern0pt{}-+++-\kern0pt{}-\kern0pt{}- & \textit{Number of cells}: 1\par \textit{Cells (identifiers)}: 7\par \textit{Multisections}: charged\_\allowbreak{}lepton\_\allowbreak{}G4\par \textit{Collections}: charged\_\allowbreak{}lepton\par \textit{Color}: none\par \textit{Joint profiles}: 1 & (60,\allowbreak{}78,\allowbreak{}25)\allowbreak{}x1 & (80,\allowbreak{}54,\allowbreak{}6)\allowbreak{}x1 \\
% VI-CATALOG-ROW familias 41
41\par 30|\allowbreak{}6|\allowbreak{}8|\allowbreak{}0|\allowbreak{}-12|\allowbreak{}+++-\kern0pt{}-\kern0pt{}-+++-\kern0pt{}-\kern0pt{}- & \textit{Number of cells}: 1\par \textit{Cells (identifiers)}: 28\par \textit{Multisections}: quark\_\allowbreak{}down\_\allowbreak{}G4\par \textit{Collections}: quark\_\allowbreak{}down\par \textit{Color}: R\par \textit{Joint profiles}: 1 & (80,\allowbreak{}54,\allowbreak{}7)\allowbreak{}x1 & (80,\allowbreak{}54,\allowbreak{}6)\allowbreak{}x1 \\
% VI-CATALOG-ROW familias 42
42\par 34|\allowbreak{}0|\allowbreak{}-12|\allowbreak{}-2|\allowbreak{}-8|\allowbreak{}0++-\kern0pt{}-\kern0pt{}-0++-\kern0pt{}-\kern0pt{}- & \textit{Number of cells}: 1\par \textit{Cells (identifiers)}: 45\par \textit{Multisections}: charged\_\allowbreak{}lepton\_\allowbreak{}G4\par \textit{Collections}: charged\_\allowbreak{}lepton\par \textit{Color}: none\par \textit{Joint profiles}: 1 & (80,\allowbreak{}66,\allowbreak{}4)\allowbreak{}x1 & (56,\allowbreak{}54,\allowbreak{}5)\allowbreak{}x1 \\
% VI-CATALOG-ROW familias 43
43\par 34|\allowbreak{}0|\allowbreak{}-12|\allowbreak{}0|\allowbreak{}-12|\allowbreak{}+++-\kern0pt{}-\kern0pt{}-+++-\kern0pt{}-\kern0pt{}- & \textit{Number of cells}: 2\par \textit{Cells (identifiers)}: 57,\allowbreak{}76\par \textit{Multisections}: charged\_\allowbreak{}lepton\_\allowbreak{}G2 | neutrino\_\allowbreak{}G4\par \textit{Collections}: charged\_\allowbreak{}lepton | neutrino\par \textit{Color}: none\par \textit{Joint profiles}: 2 & (40,\allowbreak{}84,\allowbreak{}30)\allowbreak{}x1; (80,\allowbreak{}66,\allowbreak{}2)\allowbreak{}x1 & (80,\allowbreak{}54,\allowbreak{}6)\allowbreak{}x2 \\
% VI-CATALOG-ROW familias 44
44\par 34|\allowbreak{}0|\allowbreak{}-16|\allowbreak{}0|\allowbreak{}-8|\allowbreak{}+++-0-+++-0- & \textit{Number of cells}: 1\par \textit{Cells (identifiers)}: 37\par \textit{Multisections}: charged\_\allowbreak{}lepton\_\allowbreak{}G4\par \textit{Collections}: charged\_\allowbreak{}lepton\par \textit{Color}: none\par \textit{Joint profiles}: 1 & (40,\allowbreak{}24,\allowbreak{}2)\allowbreak{}x1 & (80,\allowbreak{}54,\allowbreak{}6)\allowbreak{}x1 \\
% VI-CATALOG-ROW familias 45
45\par 34|\allowbreak{}0|\allowbreak{}-16|\allowbreak{}0|\allowbreak{}-8|\allowbreak{}+0+-\kern0pt{}-\kern0pt{}-+0+-\kern0pt{}-\kern0pt{}- & \textit{Number of cells}: 1\par \textit{Cells (identifiers)}: 15\par \textit{Multisections}: quark\_\allowbreak{}up\_\allowbreak{}G3\par \textit{Collections}: quark\_\allowbreak{}up\par \textit{Color}: B\par \textit{Joint profiles}: 1 & (40,\allowbreak{}78,\allowbreak{}27)\allowbreak{}x1 & (80,\allowbreak{}54,\allowbreak{}6)\allowbreak{}x1 \\
% VI-CATALOG-ROW familias 46
46\par 34|\allowbreak{}0|\allowbreak{}-20|\allowbreak{}-2|\allowbreak{}-8|\allowbreak{}0++-\kern0pt{}-\kern0pt{}-0++-\kern0pt{}-\kern0pt{}- & \textit{Number of cells}: 1\par \textit{Cells (identifiers)}: 35\par \textit{Multisections}: quark\_\allowbreak{}up\_\allowbreak{}G3\par \textit{Collections}: quark\_\allowbreak{}up\par \textit{Color}: B\par \textit{Joint profiles}: 1 & (100,\allowbreak{}66,\allowbreak{}0)\allowbreak{}x1 & (56,\allowbreak{}54,\allowbreak{}5)\allowbreak{}x1 \\
% VI-CATALOG-ROW familias 47
47\par 34|\allowbreak{}0|\allowbreak{}-4|\allowbreak{}-2|\allowbreak{}-8|\allowbreak{}++0-\kern0pt{}-\kern0pt{}-++0-\kern0pt{}-\kern0pt{}- & \textit{Number of cells}: 1\par \textit{Cells (identifiers)}: 56\par \textit{Multisections}: neutrino\_\allowbreak{}G3\par \textit{Collections}: neutrino\par \textit{Color}: none\par \textit{Joint profiles}: 1 & (40,\allowbreak{}24,\allowbreak{}2)\allowbreak{}x1 & (36,\allowbreak{}54,\allowbreak{}4)\allowbreak{}x1 \\
% VI-CATALOG-ROW familias 48
48\par 34|\allowbreak{}0|\allowbreak{}-4|\allowbreak{}0|\allowbreak{}-12|\allowbreak{}+++-\kern0pt{}-\kern0pt{}-+++-\kern0pt{}-\kern0pt{}- & \textit{Number of cells}: 3\par \textit{Cells (identifiers)}: 6,\allowbreak{}25,\allowbreak{}46\par \textit{Multisections}: charged\_\allowbreak{}lepton\_\allowbreak{}G2 | neutrino\_\allowbreak{}G4 | quark\_\allowbreak{}down\_\allowbreak{}G4\par \textit{Collections}: charged\_\allowbreak{}lepton | neutrino | quark\_\allowbreak{}down\par \textit{Color}: B | none\par \textit{Joint profiles}: 2 & (80,\allowbreak{}66,\allowbreak{}2)\allowbreak{}x1; (80,\allowbreak{}66,\allowbreak{}3)\allowbreak{}x2 & (80,\allowbreak{}54,\allowbreak{}6)\allowbreak{}x3 \\
% VI-CATALOG-ROW familias 49
49\par 34|\allowbreak{}0|\allowbreak{}-4|\allowbreak{}2|\allowbreak{}-8|\allowbreak{}+++-\kern0pt{}-0+++-\kern0pt{}-0 & \textit{Number of cells}: 1\par \textit{Cells (identifiers)}: 67\par \textit{Multisections}: quark\_\allowbreak{}up\_\allowbreak{}G3\par \textit{Collections}: quark\_\allowbreak{}up\par \textit{Color}: G\par \textit{Joint profiles}: 1 & (100,\allowbreak{}66,\allowbreak{}2)\allowbreak{}x1 & (56,\allowbreak{}54,\allowbreak{}5)\allowbreak{}x1 \\
% VI-CATALOG-ROW familias 50
50\par 34|\allowbreak{}0|\allowbreak{}-8|\allowbreak{}-2|\allowbreak{}-8|\allowbreak{}++0-\kern0pt{}-\kern0pt{}-++0-\kern0pt{}-\kern0pt{}- & \textit{Number of cells}: 1\par \textit{Cells (identifiers)}: 66\par \textit{Multisections}: quark\_\allowbreak{}down\_\allowbreak{}G3\par \textit{Collections}: quark\_\allowbreak{}down\par \textit{Color}: B\par \textit{Joint profiles}: 1 & (40,\allowbreak{}78,\allowbreak{}27)\allowbreak{}x1 & (80,\allowbreak{}54,\allowbreak{}6)\allowbreak{}x1 \\
% VI-CATALOG-ROW familias 51
51\par 34|\allowbreak{}0|\allowbreak{}-8|\allowbreak{}-2|\allowbreak{}-8|\allowbreak{}0++-\kern0pt{}-\kern0pt{}-0++-\kern0pt{}-\kern0pt{}- & \textit{Number of cells}: 1\par \textit{Cells (identifiers)}: 77\par \textit{Multisections}: charged\_\allowbreak{}lepton\_\allowbreak{}G4\par \textit{Collections}: charged\_\allowbreak{}lepton\par \textit{Color}: none\par \textit{Joint profiles}: 1 & (80,\allowbreak{}66,\allowbreak{}3)\allowbreak{}x1 & (56,\allowbreak{}54,\allowbreak{}5)\allowbreak{}x1 \\
% VI-CATALOG-ROW familias 52
52\par 34|\allowbreak{}0|\allowbreak{}-8|\allowbreak{}0|\allowbreak{}-12|\allowbreak{}+++-\kern0pt{}-\kern0pt{}-+++-\kern0pt{}-\kern0pt{}- & \textit{Number of cells}: 1\par \textit{Cells (identifiers)}: 36\par \textit{Multisections}: quark\_\allowbreak{}down\_\allowbreak{}G4\par \textit{Collections}: quark\_\allowbreak{}down\par \textit{Color}: G\par \textit{Joint profiles}: 1 & (40,\allowbreak{}84,\allowbreak{}30)\allowbreak{}x1 & (80,\allowbreak{}54,\allowbreak{}6)\allowbreak{}x1 \\
% VI-CATALOG-ROW familias 53
53\par 34|\allowbreak{}0|\allowbreak{}-8|\allowbreak{}0|\allowbreak{}-8|\allowbreak{}+0+-\kern0pt{}-\kern0pt{}-+0+-\kern0pt{}-\kern0pt{}- & \textit{Number of cells}: 1\par \textit{Cells (identifiers)}: 5\par \textit{Multisections}: charged\_\allowbreak{}lepton\_\allowbreak{}G4\par \textit{Collections}: charged\_\allowbreak{}lepton\par \textit{Color}: none\par \textit{Joint profiles}: 1 & (40,\allowbreak{}24,\allowbreak{}2)\allowbreak{}x1 & (80,\allowbreak{}54,\allowbreak{}6)\allowbreak{}x1 \\
% VI-CATALOG-ROW familias 54
54\par 34|\allowbreak{}0|\allowbreak{}0|\allowbreak{}-2|\allowbreak{}-8|\allowbreak{}++0-\kern0pt{}-\kern0pt{}-++0-\kern0pt{}-\kern0pt{}- & \textit{Number of cells}: 1\par \textit{Cells (identifiers)}: 26\par \textit{Multisections}: neutrino\_\allowbreak{}G3\par \textit{Collections}: neutrino\par \textit{Color}: none\par \textit{Joint profiles}: 1 & (40,\allowbreak{}78,\allowbreak{}27)\allowbreak{}x1 & (80,\allowbreak{}54,\allowbreak{}6)\allowbreak{}x1 \\
% VI-CATALOG-ROW familias 55
55\par 34|\allowbreak{}0|\allowbreak{}0|\allowbreak{}0|\allowbreak{}-8|\allowbreak{}+0+-\kern0pt{}-\kern0pt{}-+0+-\kern0pt{}-\kern0pt{}- & \textit{Number of cells}: 1\par \textit{Cells (identifiers)}: 47\par \textit{Multisections}: quark\_\allowbreak{}up\_\allowbreak{}G3\par \textit{Collections}: quark\_\allowbreak{}up\par \textit{Color}: G\par \textit{Joint profiles}: 1 & (40,\allowbreak{}78,\allowbreak{}27)\allowbreak{}x1 & (80,\allowbreak{}54,\allowbreak{}6)\allowbreak{}x1 \\
% VI-CATALOG-ROW familias 56
56\par 34|\allowbreak{}0|\allowbreak{}0|\allowbreak{}2|\allowbreak{}-8|\allowbreak{}+++0-\kern0pt{}-+++0-\kern0pt{}- & \textit{Number of cells}: 1\par \textit{Cells (identifiers)}: 16\par \textit{Multisections}: quark\_\allowbreak{}down\_\allowbreak{}G3\par \textit{Collections}: quark\_\allowbreak{}down\par \textit{Color}: G\par \textit{Joint profiles}: 1 & (40,\allowbreak{}24,\allowbreak{}2)\allowbreak{}x1 & (36,\allowbreak{}54,\allowbreak{}4)\allowbreak{}x1 \\
\end{longtable}
\endgroup
\subsection{The 23 documented classes}

These rows bring together public sectors, residences and representation.
The status key refers to the adjacent legend and preserves the status of
the source. In particular, an exact internal operator, a conditional
evaluation and a closed scalar realization remain distinct assertions.
The table neither updates their status nor turns a conditional record
into a prediction.

\begingroup
\setlength{\tabcolsep}{4pt}
\setlength{\extrarowheight}{4pt}
\renewcommand{\arraystretch}{1.13}
\begin{longtable}{@{}>{\raggedright\arraybackslash}p{\dimexpr 0.08860759\linewidth-2\tabcolsep\relax}>{\raggedright\arraybackslash}p{\dimexpr 0.41772152\linewidth-2\tabcolsep\relax}>{\raggedright\arraybackslash}p{\dimexpr 0.49367089\linewidth-2\tabcolsep\relax}@{}}
\caption{Original class statuses.}\label{tab:vi-estados-clases}\\
\toprule
Key & Literal status & Documentary reading \\
\midrule
\endfirsthead
\multicolumn{3}{l}{\textit{Continuation}}\\
\toprule
Key & Literal status & Documentary reading \\
\midrule
\endhead
\midrule
\multicolumn{3}{r}{Continued on the next page}\\
\endfoot
\bottomrule
\endlastfoot
C01 & CLOSED\_\allowbreak{}ELECTRON\_\allowbreak{}TWO\_\allowbreak{}WAY & Electron: documented bidirectional closure. \\
C02 & CLOSED\_\allowbreak{}NULL\_\allowbreak{}CHART & Null chart closed at this stage. \\
C03 & NOT\_\allowbreak{}A\_\allowbreak{}UNIVERSAL\_\allowbreak{}MASS\_\allowbreak{}SCALAR & The class is not expressed as a universal mass scalar. \\
C04 & OPEN\_\allowbreak{}FLAVOR\_\allowbreak{}COLOR\_\allowbreak{}SCALARIZATION & Flavor and color scalarization open at this stage. \\
C05 & OPEN\_\allowbreak{}MIXING\_\allowbreak{}AND\_\allowbreak{}SCALE & Mixing and scale open at this stage. \\
C06 & OPEN\_\allowbreak{}NAMED\_\allowbreak{}STATE\_\allowbreak{}SCALARIZATION & Scalarization of the individually identified state open at this stage. \\
C07 & OPEN\_\allowbreak{}PHYSICAL\_\allowbreak{}SCALARIZATION & Physical scalarization open at this stage. \\
C08 & OPEN\_\allowbreak{}ROUTE\_\allowbreak{}AND\_\allowbreak{}SCALARIZATION & Route and scalarization open at this stage. \\
\end{longtable}
\endgroup

\begingroup
\setlength{\tabcolsep}{4pt}
\setlength{\extrarowheight}{4pt}
\renewcommand{\arraystretch}{1.08}
\begin{longtable}{@{}>{\raggedright\arraybackslash}p{\dimexpr 0.17307692\linewidth-2\tabcolsep\relax}>{\raggedright\arraybackslash}p{\dimexpr 0.38461538\linewidth-2\tabcolsep\relax}>{\raggedright\arraybackslash}p{\dimexpr 0.33333333\linewidth-2\tabcolsep\relax}>{\raggedright\arraybackslash}p{\dimexpr 0.10897436\linewidth-2\tabcolsep\relax}@{}}
\caption{The 23 catalogue classes, in their documentary order.}\label{tab:vi-clases}\\
\toprule
Class & Representatives and residence & Incidence and representation & Status \\
\midrule
\endfirsthead
\multicolumn{4}{l}{\textit{Continuation}}\\
\toprule
Class & Representatives and residence & Incidence and representation & Status \\
\midrule
\endhead
\midrule
\multicolumn{4}{r}{Continued on the next page}\\
\endfoot
\bottomrule
\endlastfoot
% VI-CATALOG-ROW clases 1
SM-LC-1\par lepton\_\allowbreak{}cargado & \textit{Representatives}: e− / e+\par \textit{Residence}: center (5,\allowbreak{}5) and antisection & \textit{Selected cells}: 1\par \textit{Multisections}: charged\_\allowbreak{}lepton\_\allowbreak{}G0\par \textit{Charge}: ±1 e\par \textit{Representation}: colorless & C01 \\
% VI-CATALOG-ROW clases 2
SM-LC-2\par lepton\_\allowbreak{}cargado & \textit{Representatives}: μ− / μ+\par \textit{Residence}: charged-lepton multisection & \textit{Selected cells}: 8\par \textit{Multisections}: charged\_\allowbreak{}lepton\_\allowbreak{}G2\par \textit{Charge}: ±1 e\par \textit{Representation}: colorless & C07 \\
% VI-CATALOG-ROW clases 3
SM-LC-3\par lepton\_\allowbreak{}cargado & \textit{Representatives}: τ− / τ+\par \textit{Residence}: charged-lepton multisection & \textit{Selected cells}: 16\par \textit{Multisections}: charged\_\allowbreak{}lepton\_\allowbreak{}G4\par \textit{Charge}: ±1 e\par \textit{Representation}: colorless & C07 \\
% VI-CATALOG-ROW clases 4
SM-NU-1\par neutrino & \textit{Representatives}: ν\_\allowbreak{}e / anti-ν\_\allowbreak{}e\par \textit{Residence}: neutrino multisections & \textit{Selected cells}: 20\par \textit{Multisections}: neutrino\_\allowbreak{}G1 | neutrino\_\allowbreak{}G2 | neutrino\_\allowbreak{}G3 | neutrino\_\allowbreak{}G4\par \textit{Charge}: 0\par \textit{Representation}: colorless & C05 \\
% VI-CATALOG-ROW clases 5
SM-NU-2\par neutrino & \textit{Representatives}: ν\_\allowbreak{}μ / anti-ν\_\allowbreak{}μ\par \textit{Residence}: neutrino multisections & \textit{Selected cells}: 20\par \textit{Multisections}: neutrino\_\allowbreak{}G1 | neutrino\_\allowbreak{}G2 | neutrino\_\allowbreak{}G3 | neutrino\_\allowbreak{}G4\par \textit{Charge}: 0\par \textit{Representation}: colorless & C05 \\
% VI-CATALOG-ROW clases 6
SM-NU-3\par neutrino & \textit{Representatives}: ν\_\allowbreak{}τ / anti-ν\_\allowbreak{}τ\par \textit{Residence}: neutrino multisections & \textit{Selected cells}: 20\par \textit{Multisections}: neutrino\_\allowbreak{}G1 | neutrino\_\allowbreak{}G2 | neutrino\_\allowbreak{}G3 | neutrino\_\allowbreak{}G4\par \textit{Charge}: 0\par \textit{Representation}: colorless & C05 \\
% VI-CATALOG-ROW clases 7
SM-Q-U1\par quark & \textit{Representatives}: u / anti-u\par \textit{Residence}: up-type multisection & \textit{Selected cells}: 16\par \textit{Multisections}: quark\_\allowbreak{}up\_\allowbreak{}G1 | quark\_\allowbreak{}up\_\allowbreak{}G3\par \textit{Charge}: ±2/\allowbreak{}3 e\par \textit{Representation}: R,\allowbreak{}G,\allowbreak{}B = Z/\allowbreak{}3 orbit & C04 \\
% VI-CATALOG-ROW clases 8
SM-Q-D1\par quark & \textit{Representatives}: d / anti-d\par \textit{Residence}: down-type multisection & \textit{Selected cells}: 20\par \textit{Multisections}: quark\_\allowbreak{}down\_\allowbreak{}G1 | quark\_\allowbreak{}down\_\allowbreak{}G2 | quark\_\allowbreak{}down\_\allowbreak{}G3 | quark\_\allowbreak{}down\_\allowbreak{}G4\par \textit{Charge}: ±1/\allowbreak{}3 e\par \textit{Representation}: R,\allowbreak{}G,\allowbreak{}B = Z/\allowbreak{}3 orbit & C04 \\
% VI-CATALOG-ROW clases 9
SM-Q-U2\par quark & \textit{Representatives}: c / anti-c\par \textit{Residence}: up-type multisection & \textit{Selected cells}: 16\par \textit{Multisections}: quark\_\allowbreak{}up\_\allowbreak{}G1 | quark\_\allowbreak{}up\_\allowbreak{}G3\par \textit{Charge}: ±2/\allowbreak{}3 e\par \textit{Representation}: R,\allowbreak{}G,\allowbreak{}B = Z/\allowbreak{}3 orbit & C04 \\
% VI-CATALOG-ROW clases 10
SM-Q-D2\par quark & \textit{Representatives}: s / anti-s\par \textit{Residence}: down-type multisection & \textit{Selected cells}: 20\par \textit{Multisections}: quark\_\allowbreak{}down\_\allowbreak{}G1 | quark\_\allowbreak{}down\_\allowbreak{}G2 | quark\_\allowbreak{}down\_\allowbreak{}G3 | quark\_\allowbreak{}down\_\allowbreak{}G4\par \textit{Charge}: ±1/\allowbreak{}3 e\par \textit{Representation}: R,\allowbreak{}G,\allowbreak{}B = Z/\allowbreak{}3 orbit & C04 \\
% VI-CATALOG-ROW clases 11
SM-Q-U3\par quark & \textit{Representatives}: t / anti-t\par \textit{Residence}: up-type multisection & \textit{Selected cells}: 16\par \textit{Multisections}: quark\_\allowbreak{}up\_\allowbreak{}G1 | quark\_\allowbreak{}up\_\allowbreak{}G3\par \textit{Charge}: ±2/\allowbreak{}3 e\par \textit{Representation}: R,\allowbreak{}G,\allowbreak{}B = Z/\allowbreak{}3 orbit & C04 \\
% VI-CATALOG-ROW clases 12
SM-Q-D3\par quark & \textit{Representatives}: b / anti-b\par \textit{Residence}: down-type multisection & \textit{Selected cells}: 20\par \textit{Multisections}: quark\_\allowbreak{}down\_\allowbreak{}G1 | quark\_\allowbreak{}down\_\allowbreak{}G2 | quark\_\allowbreak{}down\_\allowbreak{}G3 | quark\_\allowbreak{}down\_\allowbreak{}G4\par \textit{Charge}: ±1/\allowbreak{}3 e\par \textit{Representation}: R,\allowbreak{}G,\allowbreak{}B = Z/\allowbreak{}3 orbit & C04 \\
% VI-CATALOG-ROW clases 13
SM-GA-EM\par boson\_\allowbreak{}gauge & \textit{Representatives}: γ\par \textit{Residence}: null gauge sector & \textit{Selected cells}: 0\par \textit{Multisections}: \textemdash{}\par \textit{Charge}: 0\par \textit{Representation}: colorless & C02 \\
% VI-CATALOG-ROW clases 14
SM-GA-GL\par boson\_\allowbreak{}gauge & \textit{Representatives}: g\par \textit{Residence}: color gauge sector & \textit{Selected cells}: 0\par \textit{Multisections}: \textemdash{}\par \textit{Charge}: 0\par \textit{Representation}: gauge octet & C02 \\
% VI-CATALOG-ROW clases 15
SM-GA-W\par boson\_\allowbreak{}gauge & \textit{Representatives}: W+ / W−\par \textit{Residence}: massive gauge sector & \textit{Selected cells}: 0\par \textit{Multisections}: \textemdash{}\par \textit{Charge}: ±1 e\par \textit{Representation}: colorless & C08 \\
% VI-CATALOG-ROW clases 16
SM-GA-Z\par boson\_\allowbreak{}gauge & \textit{Representatives}: Z0\par \textit{Residence}: massive gauge sector & \textit{Selected cells}: 0\par \textit{Multisections}: \textemdash{}\par \textit{Charge}: 0\par \textit{Representation}: colorless & C08 \\
% VI-CATALOG-ROW clases 17
SM-H-1\par escalar & \textit{Representatives}: H\par \textit{Residence}: scalar sector & \textit{Selected cells}: 0\par \textit{Multisections}: \textemdash{}\par \textit{Charge}: 0\par \textit{Representation}: colorless & C08 \\
% VI-CATALOG-ROW clases 18
COMP-MESON\par meson & \textit{Representatives}: π0, π±, K0, K±, η, ρ, φ, J/\allowbreak{}ψ, Υ…\par \textit{Residence}: product of section and antisection & \textit{Selected cells}: 0\par \textit{Multisections}: \textemdash{}\par \textit{Charge}: dependent on the composite state\par \textit{Representation}: singlet or other representations & C06 \\
% VI-CATALOG-ROW clases 19
COMP-BARYON\par barion & \textit{Representatives}: p, n and baryonic configurations\par \textit{Residence}: product of three sections & \textit{Selected cells}: 0\par \textit{Multisections}: \textemdash{}\par \textit{Charge}: dependent on the composite state\par \textit{Representation}: physical singlet & C06 \\
% VI-CATALOG-ROW clases 20
TOPO-FIB\par anyon\_\allowbreak{}no\_\allowbreak{}abeliano & \textit{Representatives}: 1, τ\par \textit{Residence}: topological sector & \textit{Selected cells}: 0\par \textit{Multisections}: \textemdash{}\par \textit{Charge}: defined by the topological codomain\par \textit{Representation}: not applicable & C03 \\
% VI-CATALOG-ROW clases 21
TOPO-ISING\par sector\_\allowbreak{}majorana & \textit{Representatives}: 1, ψ, σ\par \textit{Residence}: topological sector / zero & \textit{Selected cells}: 0\par \textit{Multisections}: \textemdash{}\par \textit{Charge}: dependent on the realization\par \textit{Representation}: colorless & C03 \\
% VI-CATALOG-ROW clases 22
TOPO-Z6\par anyones\_\allowbreak{}abelianos & \textit{Representatives}: boson, fermion and four braided sectors\par \textit{Residence}: topological sector Z/\allowbreak{}6 & \textit{Selected cells}: 0\par \textit{Multisections}: \textemdash{}\par \textit{Charge}: Z/\allowbreak{}6 definida por el trenzado\par \textit{Representation}: colorless & C03 \\
% VI-CATALOG-ROW clases 23
VIRT-1\par virtual & \textit{Representatives}: virtual sections\par \textit{Residence}: TPK inverse system & \textit{Selected cells}: 0\par \textit{Multisections}: \textemdash{}\par \textit{Charge}: dependent on the prolongation sector\par \textit{Representation}: according to sector & C03 \\
\end{longtable}
\endgroup

\input{sections/12_catalogo_datos_companero.tex}
\end{document}
